\section{Conclusion}
\label{sec:conclusion}

We built \deadeye{} to separate three things that evaluations of language-model agents usually run together: how large the model is, how it is turned into a policy, and what kind of decision it faces. \result{[Two or three sentences on the main findings: the slope under generation and how it differs between planning and one-shot tasks; whether adaptation lets a small model match a ten times larger one, and where it breaks under shift; how much of the small-model deficit scoring recovers.]} For someone choosing a model for a decision task, the practical reading is \result{[recommendation on model size and conversion method, conditional on the findings]}. The next steps follow from the limitations: reinforcement-learning conversions, learnable teachers for the bandit tasks, and longer text environments in which the oracle is harder to compute.

\paragraph{Reproducibility statement.}
The harness, configurations, prompts and analysis code are released under the Apache-2.0 license at the repository linked on the title page, and the pre-registration is frozen at the git tag \texttt{prereg-v1}. Every result table and figure is produced by \texttt{deadeye report} from logged runs, and every number in the text is taken from the report's \texttt{cells.csv} or \texttt{pairwise\_methods.csv} or from \texttt{deadeye compare}. Appendix~\ref{app:prompts} reproduces every prompt verbatim, Appendix~\ref{app:envs} the task generators, Appendix~\ref{app:prereg} the pre-registered criteria, Appendix~\ref{app:compute} the compute budget and Appendix~\ref{app:repro} the commands that rebuild the paper from a fresh clone. \todo{Archive DOI for the code, results and report.}
